% ============================================================================
% Accessibility
% ============================================================================
\DocumentMetadata{
  lang = en-US, 
  pdfstandard = ua-2, 
  tagging = on, 
  tagging-setup = {math/setup=mathml-SE}
}

%Document format
\documentclass[letterpaper,12pt]{article}

%Set table of contents depth 1=sections, 2=subsections
\setcounter{tocdepth}{2}

%Margin sizing
\usepackage[left=1in, right=1in, top=1in, bottom=1in]{geometry}

\usepackage{amsmath}
\usepackage{amssymb}

%Prevent objects from moving into a different section than where it's placed
\usepackage[section]{placeins}

%Allows pictures to be placed in document
\usepackage{graphicx}

\usepackage{booktabs}

\usepackage{multirow}

\usepackage{multicol}

\usepackage{float}

%Allows attaching of pdf documents
%\usepackage{pdfpages}

%Allows the line spacing to be changed
\usepackage{setspace}

%Package plus the option places a blank line between paragraphs
\usepackage[parfill]{parskip}

\usepackage{caption}
\usepackage{subcaption}

%allows verbatim tags
%\usepackage{verbatim}

%gives more color options for hyperlinks/bookmarks
\usepackage[table]{xcolor}

%Define colors from the Baker College Style Guide
\definecolor{BakerRed}{RGB}{239, 0, 48}%#EF0030
\definecolor{BakerADARed}{RGB}{231, 0, 46}%#E7002E
\definecolor{BakerBlack}{RGB}{31, 34, 40}%#1F2228
\definecolor{BakerWhite}{RGB}{255, 255, 255}%#FFFFFF
\definecolor{HoneyYellow}{RGB}{255, 196, 0}%#FFC400
\definecolor{UpNorthGreen}{RGB}{114, 191, 68}%#72BF44
\definecolor{LakeBlue}{RGB}{134, 183, 227}%#86B7E3
\definecolor{Abbey}{RGB}{80, 82, 88}%#505258
\definecolor{RollingStone}{RGB}{121, 124, 126}%#797C7E
\definecolor{HitGray}{RGB}{165, 169, 172}%#A5A9AC
\definecolor{Iron1}{RGB}{201, 206, 209}%#C9CED1
\definecolor{Iron}{RGB}{222, 225, 226}%#DEE1E2
\definecolor{BlackHaze1}{RGB}{237, 240, 240}%#EDF0F0
\definecolor{BlackHaze}{RGB}{245, 247, 247}%#F5F7F7

%Package for creating an index
\usepackage{imakeidx}

%Package that allows the creation of bibliographies
\usepackage[
    backend=biber,
    style=numeric-comp,
    sorting=nty,
    ]{biblatex}

%Load the bibliography file
\addbibresource{./Backmatter/reference.bib}

\usepackage[acronym]{glossaries}

\usepackage{fontspec}

\setmainfont{NHaasGroteskTXPro-}[
    Path=./NeueHaasGroteskFiles/,
    Scale=1.0,
    Extension = .ttf,
    UprightFont=*55Rg,
    BoldFont=*75Bd,
    ItalicFont=*56It,
    BoldItalicFont=*76BdIt
    ]

%Places bookmarks in the pdf file
\usepackage[luatex]{hyperref}

\hypersetup{
    bookmarks = true,
    colorlinks = true,
    linkcolor = BakerADARed,
    linkbordercolor = white,
    citecolor = BakerADARed,
    citebordercolor = white,
    urlcolor = LakeBlue,
    urlbordercolor = white,
    pdfauthor = {Brad Peirson},
    pdftitle = {Mechatronics Memorandum of Understanding},
    pdfsubject = {Mechatronics},
    pdfkeywords = {},
    pdfcreator = {LaTeX with hyperref package},
    pdfproducer = {LuaLaTeX}
    }

    % Leaving the graphics path commented out so that it can be customized per presentation
% %All of the graphics are in a separate folder, this makes it so the folder doesn't have to be repeated every time
% \graphicspath{
%     {../../assets/images}
%     {../../assets/diagrams}
%     {../../assets/plots}
% }

%Tells LaTeX to generate the files necessary for an index
%\makeindex

%Set the line spacing
%\onehalfspacing

%Adds a source element that can be used for pictures
\newcommand{\source}[1]{\vspace{-2ex} \caption{\scriptsize{Source: {#1}}} }

\usepackage[americanvoltages]{circuitikz}
\usepackage{pgfplots}
\usepackage{pgf-pie}
\usepackage{standalone}

\usetikzlibrary{
    intersections, 
    fillbetween,
    patterns,
    decorations.pathreplacing,
    decorations.pathmorphing,
    decorations.text,
    decorations.markings,
    arrows.meta,
    arrows,
    shapes,
    shapes.geometric,
    shadows,
    calligraphy,
    backgrounds,
    intersections,
    calc,
    math,
    positioning
    }

\tikzset{middlearrow/.style={
        decoration={markings, 
        mark= at position 0.5 with {\arrow{#1}},
        },
        postaction={decorate}
    }
}

\pgfplotsset{compat=1.18}

\pgfdeclarelayer{ft}
\pgfdeclarelayer{bg}
\pgfsetlayers{bg, main, ft}

% --- POWERFLEX 40 INLINE KEYPAD BUTTONS ---

% Base style for the circular buttons
\tikzset{
    pf40button/.style={
        draw=BakerBlack, 
        minimum height=1.75em,
        minimum width=3em,
        rounded corners=10pt,
        inner sep=0pt, 
        font=\bfseries,
        baseline=-0.6ex
    }
}

% --- POWERFLEX 40 INLINE KEYPAD BUTTONS ---

% Base style for the circular buttons
\tikzset{
    pf40button/.style={
        draw=BakerBlack, 
        minimum height=1.75em,
        minimum width=3em,
        rounded corners=10pt,
        inner sep=0pt, 
        font=\bfseries,
        baseline=-0.6ex
    }
}

% 1. ESCAPE BUTTON
\newcommand{\pfEsc}{%
    \begin{tikzpicture}[alt={PowerFlex 40 Escape button. A rounded blue rectangle with white Esc inside.}]
        \node[pf40button, fill=LakeBlue, text=white] {Esc};
    \end{tikzpicture}%
}

% 2. SELECT BUTTON
\newcommand{\pfSel}{%
    \begin{tikzpicture}[alt={PowerFlex 40 Select button. A rounded blue rectangle with white Sel inside.}]
        \node[pf40button, fill=LakeBlue, text=white] {Sel};
    \end{tikzpicture}%
}

% 3. ENTER BUTTON
\newcommand{\pfEnter}{%
    \begin{tikzpicture}[alt={PowerFlex 40 Enter button. A rounded blue rectangle with white arrow inside.}]
        % Draw the base gray circle button
        \node(btn)[pf40button, fill=LakeBlue] {};
        
        % Draw the carriage return / elbow arrow inside the button
        \draw[-{Stealth[length=8pt, width=8pt]}, white, line width=3pt] 
            (btn.center) ++(0.65em, 0.4em) -- ++(0,-0.5em) -- ++(-1.2em,0);;
    \end{tikzpicture}%
}

% 4. UP ARROW BUTTON
\newcommand{\pfUp}{%
    \begin{tikzpicture}[alt={PowerFlex 40 Up button. A rounded blue rectangle with white triangle inside. The triangle is oriented with one point up.}]
        \node(btn)[pf40button, fill=LakeBlue] {};
        \node(tri)[draw=white, fill=white, regular polygon, regular polygon sides=3,minimum size=1em, inner sep=0, outer sep=0]{};
    \end{tikzpicture}%
}

% 5. DOWN ARROW BUTTON
\newcommand{\pfDown}{%
    \begin{tikzpicture}[alt={PowerFlex 40 Down button. A rounded blue rectangle with white triangle inside. The triangle is oriented with one point down.}]
        \node(btn)[pf40button, fill=LakeBlue] {};
        \node(tri)[draw=white, fill=white, regular polygon, regular polygon sides=3,minimum size=1em, inner sep=0, outer sep=0,rotate=180]{};
    \end{tikzpicture}%
}

% 6. START BUTTON (Green Circle with White Triangle)
\newcommand{\pfStart}{%
    \begin{tikzpicture}[alt={PowerFlex 40 Up button. A rounded green rectangle with a thick, vertical white line inside.}]
        \node(btn)[pf40button, fill=UpNorthGreen, draw=BakerBlack] {};
        \draw[white, line width=4pt](btn.center) ++(0,0.5em) -- ++(0,-1em);
    \end{tikzpicture}%
}

% 7. REVERSE BUTTON (Blue Circle with White Directional Arrow)
\newcommand{\pfRev}{%
    \begin{tikzpicture}[alt={PowerFlex 40 Reverse button. A rounded grey rectangle with an arc inside. The arc has arrows on each end.}]
        \node(btn)[pf40button, fill=blue!60!black, draw=blue!40!black] {};
        % Recreates the small hooked rotation arrow icon
        \draw[white, line width=0.8pt, ->, >=stealth] (btn.center) ++(-0.2em, 0.1em) arc (135:-135:0.2em and 0.15em);
    \end{tikzpicture}%
}

% 8. STOP BUTTON (Red Circle with White Square)
\newcommand{\pfStop}{%
    \begin{tikzpicture}[alt={PowerFlex 40 Stop button. A rounded red rectangle with white circle inside.}]
        \node(btn)[pf40button, fill=BakerADARed, draw=BakerBlack, rounded corners=6pt] {}; % slightly larger like the real VFD
        \node[draw=white, ultra thick, circle, minimum size=1em, inner sep=0, outer sep=0] at (btn.center) {};
    \end{tikzpicture}%
}